\chapter{Conclusions}
\label{chap:conclusion}

\section{What the work establishes}
\label{sec:conclusion-established}

Three claims are supported by the evidence in Chapters~\ref{chap:validation}
and~\ref{chap:results}.

The first is numerical. A refined lattice kinetic hierarchy can be made second
order accurate across refinement interfaces, and the requirement is a rescaling
of the non equilibrium population rather than a better interpolation of the
populations themselves. Table~\ref{tab:interface-error} shows that without the
rescaling the interface error is independent of the grid, so the failure is not
one that a finer mesh repairs. This resolves a question that has limited refined
lattice solvers at the Reynolds numbers where adaptivity is worth having.

The second is methodological. Refining on the resolved fraction of the Ozmidov
scale rather than on a velocity gradient produces a hierarchy that tracks the
active mixing regions, and it does so at a cost that does not grow with
stratification. Table~\ref{tab:case-matrix} holds the saving at between twenty
three and twenty seven times across the whole matrix. A gradient based criterion
applied to the same cases refines the strongly sheared but sterile layers and
achieves an active mixing fraction of $0.71$ at case S09, where the Ozmidov
criterion achieves $0.958$.

The third is physical. The turbulent Prandtl number in a stably stratified
channel is not a constant, does not depend on the gradient Richardson number
alone, and its apparent saturation near $2.0$ in the literature is a resolution
artefact rather than a property of the flow. The closure of
Equation~\eqref{eq:closure} carries the resolved fraction explicitly and reduces
the error in predicted scalar flux from 34 percent to 9 percent.

\section{Limitations}
\label{sec:conclusion-limits}

The campaign is channel flow. A channel has a fixed forcing, no large scale
pressure gradient variation, and no surface heterogeneity, and all three matter
in the atmospheric boundary layer that motivates the work. The closure
of~\eqref{eq:closure} should be read as calibrated for this configuration and
tested elsewhere before being trusted.

The Reynolds number is 950, which is high for a resolved stratified simulation
and low for the atmosphere. The buoyancy Reynolds numbers in
Table~\ref{tab:case-matrix} fall below 100 for the four most stable cases, which
is the regime where the flow is genuinely intermittent, so the physical regime of
interest is reached. Whether the closure coefficients drift with $\Rey_\tau$ is
not settled by this work, and the two point check at $\Rey_\tau = 590$ reported
in Section~\ref{sec:val-channel} constrains only the neutral end.

Case S11 is resolution limited by the solver's own diagnostic. It is reported
because the collapse behaviour is informative, and it is excluded from every
quantitative fit.

The scalar is a single active scalar with a linear background gradient. Two
scalar problems, where heat and a solute have different diffusivities and
compete for the same overturning events, are not addressed and are known to
behave differently near the collapse threshold \citep{nakagawa2018twoscalar}.

\section{Future work}
\label{sec:conclusion-future}

Three directions follow directly.

The load imbalance transient described in Section~\ref{sec:results-cost} is a
scheduling problem rather than a physics problem. The rebalance currently
assigns work by static weight $8^{\ell}$, which is correct in aggregate and
wrong for the first hundred steps after a hierarchy change because newly refined
blocks carry a larger share of interface work. A measured weight, updated from
the previous interval, should remove most of it.

The refinement criterion uses a block averaged dissipation, which smooths across
the thin sheets that the criterion exists to capture. A criterion evaluated on
the non equilibrium stress at each site, with a block level reduction that takes
a high quantile rather than a mean, would refine on the sheet rather than on the
block that contains it. The cost is one additional reduction per block per
regridding pass, which Table~\ref{tab:cost-breakdown} suggests is affordable.

Finally, the closure of~\eqref{eq:closure} carries a resolved fraction that a
large eddy simulation does not have access to, because it does not know its own
dissipation rate independently of its subgrid model. Making the closure usable
in that setting requires an estimator for $r$ from resolved quantities alone.
The scalar variance spectra of Figure~\ref{fig:spectra} suggest that the ratio
of variance above and below a fixed wavenumber is a usable proxy
\citep{devereux2023les}, and testing that is the natural next step.
